\bookchapter{Bab 48: Mengenal Android — Dari Sistem Operasi hingga Ekosistem Pengembangan}

\chapterintro{Selamat atas keputusan ini—dari “orang yang memakai ponsel” menjadi “orang yang membuat aplikasi ponsel”. Sebelum mulai praktik, mari kita buka petanya dulu: dari mana Android berasal? Mengapa pengembang di seluruh dunia rajin membuat aplikasi untuknya? Apa saja yang terjadi dari jari menyentuh keyboard sampai aplikasinya ada di ponsel pengguna?}

\booksection{Silsilah Android: dari “robot kecil” menjadi “sistem terbesar di dunia”}

Kisah Android dimulai pada 2003—seorang insinyur bernama Andy Rubin mendirikan perusahaan Android dengan tujuan yang sederhana: membuat sistem operasi pintar yang “lebih paham kamera”. Kemudian Google mengakuisisinya pada 2005, pada 2007 bersama puluhan vendor mendirikan Open Handset Alliance (OHA), dan pada September 2008 ponsel Android pertama, HTC Dream (G1), resmi dijual.

Cerita setelahnya sudah kamu tahu: sampai 2026, Android adalah sistem operasi mobile dengan pangsa pasar terbesar di dunia, dengan jumlah perangkat aktif dihitung dalam miliaran. Ponselmu, kotak TV, layar besar di mobil, jam pintarmu—kemungkinan besar semuanya menjalankan Android. Apa artinya? Artinya setiap baris kode yang kamu tulis berpeluang dipakai oleh paling banyak orang di dunia.

\begin{tcbinfo}{Tips: gen open source Android}
Kode inti Android (AOSP, Android Open Source Project) bersifat open source—vendor mana pun bebas mengambilnya, memodifikasinya, dan menyesuaikannya. Karena itu ada MIUI milik Xiaomi (sekarang bernama HyperOS), lapisan kompatibilitas HarmonyOS milik Huawei, dan One UI milik Samsung. Ini kabar baik bagi pengembang: yang kamu pelajari adalah “pengembangan Android”, bukan “pengembangan merek tertentu”—satu keahlian, berlaku di semua merek.

\end{tcbinfo}

Evolusi nomor versinya juga menarik: mulai dari Cupcake (kue cup) di versi 1.5, setiap versi besar dinamai dengan nama makanan penutup—Donut, Eclair, Frozen Yogurt, Gingerbread... sampai Android 10 pada 2019 akhirnya berganti menjadi angka murni. Per Agustus 2026, versi resmi terbaru adalah Android 17, sementara “versi dukungan minimum” yang benar-benar perlu kamu pedulikan saat mengembangkan biasanya dipatok aplikasi populer antara Android 8 (API 26) dan Android 10 (API 29)—mendukung perangkat beberapa tahun lalu membuat cakupan penggunamu cukup besar.

\begin{center}
\begin{tabularx}{\linewidth}{|X|X|X|}
\hline
\textbf{Versi} & \textbf{Kode nama/tahun} & \textbf{Hal penting bagi pengembang} \tabularnewline
\hline
Android 8 (API 26) & Oreo · 2017 & Kanal notifikasi, pembatasan latar belakang—batas minimum banyak perangkat lama \tabularnewline
Android 10 (API 29) & 2019 & Munculnya penyimpanan terpartisi (scoped storage), aturan akses berkas berubah besar \tabularnewline
Android 12 (API 31) & 2021 & Material You dengan warna dinamis, tampilan antarmuka terdongkrak besar \tabularnewline
Android 13 (API 33) & 2022 & Izin notifikasi saat runtime, dialog bertanya "boleh mengirim notifikasi atau tidak" \tabularnewline
Android 14-15 (API 34-35) & 2023-2024 & Gestur kembali prediktif, halaman 16KB, privasi makin ketat \tabularnewline
Android 16-17 (API 36-37) & 2025-2026 & Ritme versi minor lebih cepat, integrasi kemampuan AI (seperti ekstensi Gemini) \tabularnewline
\hline
\end{tabularx}
\end{center}

\begin{center}\small\color{gray}Sejarah singkat versi Android (per Agustus 2026, hal yang paling diperhatikan saat pengembangan)\end{center}

\booksection{Menulis dengan bahasa apa: Java atau Kotlin?}

Dalam sejarah pengembangan Android ada dua bahasa utama. Java adalah “bapak pendiri”—dipakai sejak Android pertama lahir; Kotlin adalah “warga kelas satu” yang ditunjuk Google pada 2017, dan sejak 2019 posisi resminya “Kotlin dulu” (Kotlin-first). Menjelang 2026, hampir semua proyek baru memakai Kotlin, dan proyek-proyek lama pun sebagian besar sedang bermigrasi.

Kata jujur saja: kalau kamu benar-benar tanpa dasar, langsung belajar Kotlin, jangan mundur ke Java. Ada tiga alasan: pertama, jumlah kode Kotlin 30\%-40\% lebih sedikit daripada Java (fungsi yang sama, Kotlin lebih pendek dan jernih); kedua, Kotlin secara bawaan mencegah “null pointer”, mimpi buruk para pemula; ketiga, semua tutorial modern, dokumentasi resmi, dan alat pemrograman AI (masih ingat buku pemrograman AI kita?) menjadikan Kotlin sebagai standar. Java hanya akan kamu temui saat merawat proyek lama atau membaca kode lama.

Gaya Java (versi bertele-tele)

\begin{lstlisting}
public class MainActivity extends Activity \{
@Override
protected void onCreate(Bundle b) \{
super.onCreate(b);
setContentView(R.layout.activity);
\}
\}
\end{lstlisting}

Gaya Kotlin (versi ringkas)

\begin{lstlisting}
class MainActivity : ComponentActivity() \{
override fun onCreate(b: Bundle?) \{
super.onCreate(b)
setContent \{ App() \}
\}
\}
\end{lstlisting}

Untuk hal yang sama, Kotlin menghemat banyak kurung kurawal dan tulisan boilerplate. Bab 50 nanti kita pelajari resmi; untuk sekarang cukup ingat: “bahasa nasional” pengembangan Android bernama Kotlin.

\booksection{Bagaimana sebuah App lahir?}

Banyak orang mengira “mengembangkan App” itu membuka komputer lalu mengetik kode sekuat hati. Nyatanya, dari ide sampai berada di ponsel pengguna, sebuah App menempuh lini produksi kurang lebih seperti ini:

Ide

Desain (tampilan+fungsi)

Menulis kode

Kompilasi dan kemas (APK/AAB)

Pengujian (emulator+perangkat nyata)

Rilis (store/paket instalasi)

\begin{center}\small\color{gray}Gambar 1: Alur lengkap sebuah App Android dari ide hingga rilis di toko\end{center}

Ada dua konsep dalam alur ini yang harus kamu kenal sekarang:

\begin{itemize}
  \item APK: format paket instalasi aplikasi Android (.apk), bisa langsung dipasang ke ponsel dengan sekali klik, seperti .exe di Windows. Untuk pengujian sendiri atau sekadar membagikan ke teman, pakailah ini.
  \item AAB: Android App Bundle (.aab), "format pengemasan premium" untuk rilis di Google Play. Toko akan memotongnya secara dinamis sesuai model ponsel pengguna sehingga paket yang diunduh tiap orang lebih kecil. Tidak bisa dipasang sendiri, hanya bisa diserahkan ke toko.
\end{itemize}

Selain itu kamu akan terus mendengar satu istilah: SDK (Software Development Kit, perangkat lunak pengembangan). Inilah satu set “perlengkapan pengembangan” lengkap dari Android resmi—mencakup pustaka API berbagai versi, alat debugging, citra emulator, dan sebagainya. Kelas-kelas seperti Activity dan Button yang kamu panggil saat menulis kode semuanya berasal dari SDK.

\booksection{Peta penuh tumpukan teknologi Android: “daftar senjata”-mu}

Sebelum belajar resmi, kenali dulu konfigurasi standar proyek Android populer tahun 2026—tidak masalah kalau daftar ini belum kamu pahami sekarang; kembali lagi setelah selesai membaca buku ini, semuanya akan terasa jernih:

\begin{center}
\begin{tabularx}{\linewidth}{|X|X|X|}
\hline
\textbf{Lapisan} & \textbf{Pilihan populer (2026)} & \textbf{Untuk apa} \tabularnewline
\hline
Bahasa & Kotlin 2.x & Bahasa untuk menulis logika \tabularnewline
Antarmuka & Jetpack Compose + Material 3 & Menggambar antarmuka lewat kode (jalur utama buku ini) \tabularnewline
Arsitektur & MVVM + ViewModel & Memisahkan antarmuka dan logika agar kode tidak berantakan \tabularnewline
Penyimpanan data & Room (database) + DataStore (konfigurasi) & Menyimpan data ke dalam ponsel \tabularnewline
Jaringan & Retrofit + OkHttp + coroutine Kotlin & Berbicara dengan server \tabularnewline
Alat build & Gradle + Kotlin DSL & Mengubah kode menjadi paket instalasi \tabularnewline
Lingkungan pengembangan & Android Studio & IDE "paket lengkap" resmi Google \tabularnewline
\hline
\end{tabularx}
\end{center}

\begin{center}\small\color{gray}Sekilas tumpukan teknologi Android populer tahun 2026\end{center}

Jangan gentar oleh rangkaian istilah berbahasa Inggris ini. Semuanya seperti panci dan wajan di dapur—untuk sekarang kamu cukup tahu “memasak harus masuk dapur”; cara memakai masing-masing akan kita bedah satu per satu di bab-bab berikutnya. Kabar baiknya: saat membuat proyek baru, Android Studio otomatis menyiapkan seluruh rangkaian ini untukmu, dan kamu tinggal memasak.

\booksection{Android vs iOS: “pelajaran memilih kubu” pertama bagi pengembang}

Mungkin kamu sedang memegang perangkat Android dan Apple sekaligus. Sebagai pengembang, perbedaan pengalaman kedua sisi patut dipahami sejak awal—bukan menentukan sikap “mana yang lebih baik”, melainkan menghitung secara realistis “sisi mana yang lebih cocok untukmu memulai”:

\begin{center}
\begin{tabularx}{\linewidth}{|X|X|X|}
\hline
\textbf{Aspek} & \textbf{Pengembangan Android} & \textbf{Pengembangan iOS} \tabularnewline
\hline
Komputer pengembangan & Windows / Mac / Linux semuanya bisa & Wajib Mac (Xcode eksklusif) \tabularnewline
Alat pengembangan & Android Studio (gratis) & Xcode (gratis) + biaya tahunan rilis 99 dolar AS \tabularnewline
Tinjauan rilis di toko & Longgar, beberapa jam hingga beberapa hari & Ketat, 1-7 hari adalah hal biasa \tabularnewline
Pasang dan coba di perangkat & APK bisa dibagikan dan dipasang dengan bebas & Distribusi dibatasi; pengujian butuh penandatanganan/TestFlight \tabularnewline
Lingkungan pendapatan & Terutama iklan/kanal lokal & Kebiasaan membayar pembelian dalam aplikasi (in-app purchase) di App Store lebih baik \tabularnewline
Biaya belajar & Konsepnya lebih banyak (empat komponen utama, dll.) & Ekosistem tertutup tapi sistemnya seragam, dokumentasinya prima \tabularnewline
\hline
\end{tabularx}
\end{center}

\begin{center}\small\color{gray}Perbandingan jalur pengembangan Android vs iOS (Agustus 2026)\end{center}

Jujur saja: kalau komputermu bukan Mac, soal ini tidak perlu dipikirkan lagi—Android adalah satu-satunya pilihan. Kalau kamu punya Mac, Bagian 8 dan Bagian 9 buku ini bisa dibaca berdampingan karena pola pikirnya saling terkait; setelah menguasai salah satunya, mempelajari yang lain akan jauh lebih cepat.

\begin{tcbpractice}{Praktik: Lihat 'Android' di Dalam Ponselmu}
Buka [Pengaturan] [Tentang Ponsel] (atau 'Tentang Perangkat') di ponselmu, cari baris 'Versi Android'—ketuk baris itu beberapa kali berturut-turut, kamu bahkan bisa memunculkan permainan easter egg bawaan sistem. Catat nomor versimu, lalu cocokkan dengan tabel versi di atas untuk melihat generasinya. Ponselmu adalah 'perangkat uji' pertamamu.

\end{tcbpractice}

\begin{tcbsummary}{Ringkasan Bab Ini}
\begin{itemize}
  \item Android lahir tahun 2008; pada 2026 ia adalah sistem mobile dengan pangsa pasar terbesar di dunia, dengan perangkat aktif mencapai miliaran
  \item Versi resmi terbaru Android 17; aplikasi populer mendukung minimal Android 8-10 agar cakupan penggunanya cukup besar
  \item Proyek baru selalu pakai Kotlin: kode hemat 30\%+, bebas null pointer, pilihan resmi Google
  \item APK dipasang sendiri, AAB untuk toko aplikasi; SDK adalah paket peralatan pengembangan resmi
  \item Tumpukan teknologi standar 2026: Kotlin + Compose + MVVM + Room + Retrofit + Gradle + Android Studio
\end{itemize}

\end{tcbsummary}
